\documentclass{article}
\usepackage{enumitem}
\usepackage{amsmath}
\begin{document}

\title{Chapter 15: Plant Growth and Development}
\date{}
\maketitle

\begin{enumerate}[label=Q\arabic*.]

\item Which growth curve best describes plant growth under normal conditions?
\\(A) Linear (arithmetic) growth
\\(B) Sigmoid (S-shaped) curve with lag, exponential, and stationary phases
\\(C) Exponential growth without limit
\\(D) Declining curve from germination

\item Auxin (IAA) promotes cell elongation by:
\\(A) Directly increasing turgor pressure
\\(B) Activating proton pumps that acidify the cell wall, activating expansins that loosen wall
\\(C) Stimulating cell division at apical meristems
\\(D) Increasing chlorophyll synthesis

\item The phenomenon of apical dominance is caused by:
\\(A) Gibberellins produced in roots
\\(B) Auxin produced at the shoot apex that inhibits lateral bud growth
\\(C) Cytokinin produced in leaves
\\(D) Ethylene produced in seeds

\item Gibberellins are primarily responsible for:
\\(A) Root elongation and gravitropism
\\(B) Stem elongation, breaking dormancy, and promoting germination of cereal seeds
\\(C) Leaf senescence and abscission
\\(D) Fruit ripening and abscission

\item Assertion: Cytokinins promote cell division and delay senescence in leaves.\\
Reason: Cytokinins stimulate synthesis of proteins and retard protein breakdown in leaves.
\\(A) Both true, Reason is correct explanation
\\(B) Both true, Reason is not correct explanation
\\(C) Assertion true, Reason false
\\(D) Assertion false, Reason true

\item Abscisic acid (ABA) is called the ``stress hormone'' because:
\\(A) It promotes rapid growth during stress
\\(B) It promotes stomatal closure during water stress and induces seed dormancy
\\(C) It stimulates ethylene production during stress
\\(D) It inhibits photosynthesis under high light stress

\item Ethylene is a unique plant hormone because:
\\(A) It is the only gaseous plant hormone
\\(B) It promotes cell elongation in all tissues
\\(C) It is produced only in roots
\\(D) It promotes seed germination exclusively

\item Phototropism in plants is caused by:
\\(A) Unequal distribution of gibberellins on lit and shaded sides
\\(B) Lateral redistribution of auxin, with more auxin on the shaded side causing differential growth
\\(C) Differential distribution of cytokinins
\\(D) Stomatal aperture differences on either side

\item Vernalization refers to:
\\(A) Conversion of vegetative growth to reproductive growth by temperature
\\(B) Promotion of flowering in plants by exposure to prolonged cold (winter) temperature
\\(C) Seed dormancy induced by cold temperature
\\(D) Dormancy breaking by cold treatment in seeds

\item The critical day length in photoperiodism determines:
\\(A) The rate of photosynthesis
\\(B) Whether a plant will flower based on the duration of darkness relative to a threshold
\\(C) The length of the growing season
\\(D) The rate of transpiration

\item In short-day plants (SDPs), flowering is promoted by:
\\(A) Short day length (long light period)
\\(B) Long uninterrupted dark period exceeding the critical night length
\\(C) Equal day and night length
\\(D) Brief red light interruption of the night period

\item Phytochrome Pfr (active form) is converted back to Pr (inactive form) by:
\\(A) Red light (660 nm)
\\(B) Far-red light (730 nm) or in darkness
\\(C) Blue light (450 nm)
\\(D) UV light

\item Which of the following is the correct sequence of events in seed germination?
\\(A) Imbibition $\rightarrow$ Activation of enzymes $\rightarrow$ Radicle emergence $\rightarrow$ Plumule growth
\\(B) Radicle emergence $\rightarrow$ Imbibition $\rightarrow$ Enzyme activation
\\(C) Enzyme activation $\rightarrow$ Imbibition $\rightarrow$ Radicle emergence
\\(D) Plumule growth $\rightarrow$ Radicle emergence $\rightarrow$ Imbibition

\item Match the plant hormones with their commercial applications:\\
1. Auxin $\rightarrow$ (a) Rooting of cuttings\\
2. Gibberellin $\rightarrow$ (b) Increasing sugarcane yield by stem elongation\\
3. Ethylene $\rightarrow$ (c) Uniform fruit ripening for market\\
4. Cytokinin $\rightarrow$ (d) Delaying senescence in stored vegetables
\\(A) 1-a, 2-b, 3-c, 4-d
\\(B) 1-b, 2-a, 3-d, 4-c
\\(C) 1-c, 2-d, 3-a, 4-b
\\(D) 1-d, 2-c, 3-b, 4-a

\item Growth inhibitors in seeds are responsible for:
\\(A) Promoting precocious germination
\\(B) Maintaining seed dormancy until favorable conditions arrive
\\(C) Accelerating embryo development
\\(D) Stimulating root growth in germinating seeds

\end{enumerate}
\end{document}
